\documentclass[12pt,a4paper]{article}

\usepackage[margin=2.5cm]{geometry}
\usepackage{fontspec}
\usepackage{polyglossia}
\setmainlanguage{vietnamese}
\setmainfont{Times New Roman}

\usepackage{enumitem}
\usepackage{longtable}
\usepackage{hyperref}
\usepackage{xcolor}
\usepackage{fancyhdr}
\usepackage{titlesec}
\usepackage{parskip}
\usepackage{booktabs}
\usepackage{caption}
\usepackage{float}
\usepackage{listings}

\hypersetup{
  colorlinks=true,
  linkcolor=black,
  urlcolor=blue,
  citecolor=black
}

\pagestyle{fancy}
\fancyhf{}
\fancyhead[L]{Báo cáo tiến độ — DevSecOps + GitOps cho Microservices trên AWS}
\fancyhead[R]{Hồ Nhật Minh — 23520924}
\fancyfoot[C]{\thepage}
\renewcommand{\headrulewidth}{0.4pt}

\titleformat{\section}{\large\bfseries}{}{0em}{}[\titlerule]
\titleformat{\subsection}{\normalsize\bfseries}{}{0em}{}

\setlength{\parskip}{6pt}
\setlength{\parindent}{0pt}
\setlength{\headheight}{15.0434pt}
\addtolength{\topmargin}{-3.0434pt}

\lstset{
  basicstyle=\ttfamily\small,
  breaklines=true,
  frame=single,
  backgroundcolor=\color{gray!5},
  columns=flexible
}

% ── Document ──────────────────────────────────────────

\begin{document}

% --- Title page ---
\begin{center}
  {\Large \textbf{ĐẠI HỌC QUỐC GIA TP. HỒ CHÍ MINH}}\\[2pt]
  {\large \textbf{TRƯỜNG ĐẠI HỌC CÔNG NGHỆ THÔNG TIN}}\\[2pt]
  {\large KHOA MẠNG MÁY TÍNH \& TRUYỀN THÔNG}\\[20pt]

  \rule{\textwidth}{1pt}\\[16pt]

  {\LARGE \textbf{BÁO CÁO TIẾN ĐỘ ĐỒ ÁN CHUYÊN NGÀNH}}\\[10pt]

  {\Large Thiết kế và Triển khai Quy trình DevSecOps}\\[4pt]
  {\Large Kết hợp GitOps cho Hệ thống Microservices trên Nền tảng AWS}\\[16pt]

  \rule{\textwidth}{1pt}\\[20pt]

  \begin{tabular}{rl}
    \textbf{Giảng viên hướng dẫn:} & Lê Anh Tuấn \\[6pt]
    \textbf{Sinh viên thực hiện:}  & Hồ Nhật Minh \\[6pt]
    \textbf{MSSV:}                 & 23520924 \\[6pt]
  \end{tabular}

  \vfill
  {\normalsize TP. Hồ Chí Minh}\\[4pt]
\end{center}

\newpage

% --- Table of Contents ---
\tableofcontents
\newpage

% ════════════════════════════════════════════════════════
\section{Tổng quan đề tài}

\subsection{Mục tiêu}

Đề tài hướng tới thiết kế và triển khai một quy trình DevSecOps tích hợp GitOps để tự động hóa việc build, kiểm thử bảo mật và triển khai hệ thống microservices lên nền tảng AWS. Quy trình đảm bảo:

\begin{itemize}[leftmargin=*,nosep]
  \item Mọi thay đổi mã nguồn được kiểm thử và quét bảo mật trước khi đến môi trường production.
  \item Việc triển khai tuân theo mô hình GitOps — toàn bộ trạng thái hệ thống được khai báo trong Git, Argo CD tự động đồng bộ lên Kubernetes.
  \item Bảo mật được tích hợp từ đầu (shift-left security): SAST, dependency scan, container image scan.
  \item Hạ tầng được quản lý 100\% bằng Infrastructure as Code (Terraform).
\end{itemize}

\subsection{Phạm vi công nghệ}

\begin{table}[H]
\centering
\begin{tabular}{@{}ll@{}}
\toprule
\textbf{Thành phần} & \textbf{Công nghệ} \\
\midrule
Hạ tầng (IaC)       & Terraform \textgreater\ 1.5, AWS Provider \textgreater\ 5.0 \\
Kubernetes           & Amazon EKS v1.33, managed node groups \\
CI/CD                & Jenkins (CI) + Argo CD (CD GitOps) \\
Bảo mật              & SonarQube SAST, OWASP Dependency Check, Trivy \\
Container Registry   & Amazon ECR \\
Giám sát             & CloudWatch + Prometheus + Grafana \\
Sample Application   & aws-containers/retail-store-sample-app (5 microservices) \\
\bottomrule
\end{tabular}
\end{table}

\subsection{Sample Application}

Hệ thống mẫu là ứng dụng bán lẻ trực tuyến gồm 5 microservice:

\begin{table}[H]
\centering
\begin{tabular}{@{}llll@{}}
\toprule
\textbf{Service}    & \textbf{Ngôn ngữ} & \textbf{Framework}       & \textbf{Backend} \\
\midrule
UI                  & Java 21           & Spring Boot + Thymeleaf  & WebFlux API Gateway \\
Cart                & Java 21           & Spring Boot              & DynamoDB \\
Orders              & Java 21           & Spring Boot              & DynamoDB, SQS, SNS, Secrets Manager \\
Catalog             & Go                & Gin                      & MySQL/MariaDB \\
Checkout            & TypeScript        & NestJS 11                & Redis \\
\bottomrule
\end{tabular}
\end{table}

% ════════════════════════════════════════════════════════
\section{Kiến trúc hệ thống}

\subsection{Sơ đồ tổng thể}

\begin{verbatim}
   Developer
      |
      v
   GitHub App Repo
      |
      v
   Jenkins CI  ----+---- SonarQube SAST
                   |---- OWASP Dependency Check
                   |---- Docker Build
                   |
                   v
              Trivy Scan
                   |
                   v (push image)
              Amazon ECR
                   |
                   v (update tag)
         GitHub GitOps Repo
                   |
                   v (pull & sync)
           Argo CD ----> EKS Cluster
                     (uit-devsecops-dev)
\end{verbatim}

\subsection{Hạ tầng AWS đã triển khai}

\begin{verbatim}
AWS Region: ap-southeast-1 (Singapore)
---------------------------------------------------
VPC (10.0.0.0/16)
  Public Subnets x 3 (1 per AZ)
    -> Internet Gateway (IGW)
    -> Application Load Balancer (ALB, public-facing)
  Private Subnets x 3 (1 per AZ)
    -> EKS Worker Nodes x 2 (t3.medium)
    -> NAT Gateway (single, shared)

EKS Control Plane v1.33
  -> VPC-CNI, CoreDNS, kube-proxy
  -> AWS Load Balancer Controller
  -> Cert Manager
  -> OTEL Operator (ADOT)

Supporting services:
  -> KMS (secrets encryption)
  -> OIDC Provider (IRSA)
  -> CloudWatch Logs
---------------------------------------------------
S3:  uit-devsecops-gitops-dev-xxxxx-tfstate (state)
DDB: uit-devsecops-gitops-dev-xxxxx-tflock (lock)
\end{verbatim}

\subsection{Mô hình 2-repo GitOps}

Hệ thống sử dụng chiến lược tách biệt hai repository:

\begin{enumerate}[leftmargin=*,nosep]
  \item \textbf{App Repo} — chứa mã nguồn microservices, Dockerfile, Jenkinsfile. Developer push code vào đây để kích hoạt CI pipeline.
  \item \textbf{GitOps Repo} — chứa các manifest triển khai (Kustomize/Helm). CI pipeline tự động cập nhật image tag vào repo này sau khi build thành công. Argo CD theo dõi repo này và tự động đồng bộ lên EKS.
\end{enumerate}

Lợi ích của mô hình này:
\begin{itemize}[leftmargin=*,nosep]
  \item Tách biệt quyền: developer chỉ cần quyền trên App Repo, không cần quyền trực tiếp lên production.
  \item Audit trail hoàn chỉnh: mọi thay đổi production đều có commit tương ứng trong GitOps Repo.
  \item Rollback = git revert: chỉ cần revert commit trong GitOps Repo, Argo CD tự động quay về trạng thái cũ.
\end{itemize}

% ════════════════════════════════════════════════════════
\section{Khối lượng đã hoàn thành}

\subsection{Giai đoạn 1: Lập kế hoạch và thiết lập}

\begin{longtable}{p{0.32\textwidth} p{0.64\textwidth}}
\toprule
\textbf{Hạng mục} & \textbf{Nội dung đã thực hiện} \\
\midrule
Cập nhật đề tài &
Thống nhất tên đề tài mới: \textit{Thiết kế và Triển khai Quy trình DevSecOps kết hợp GitOps cho Hệ thống Microservices trên Nền tảng AWS}. \\
\midrule
So sánh mô hình cũ/mới &
Đã phân tích chi tiết sự khác biệt giữa mô hình cũ (Java/Tomcat triển khai thủ công lên EC2) và mô hình mới (Microservices/EKS/GitOps). Bảng so sánh gồm 8 tiêu chí: deploy strategy, rollback, security scanning, audit trail, infrastructure management, scaling, secrets management, observability. \\
\midrule
Chọn sample app &
Đã đánh giá 3 ứng dụng mẫu do GVHD đề xuất và chọn \texttt{aws-containers/retail-store-sample-app} vì: đa dạng ngôn ngữ (Java, Go, TypeScript), kiến trúc microservices đầy đủ, hỗ trợ DynamoDB/Redis, có sẵn Dockerfile và Terraform module. \\
\midrule
Sơ đồ kiến trúc &
Đã vẽ sơ đồ tổng thể quy trình DevSecOps + GitOps bằng Mermaid, bao gồm toàn bộ luồng từ Developer → CI → Security Gates → ECR → GitOps → Argo CD → EKS. \\
\midrule
Workspace đồ án &
Đã tạo cấu trúc thư mục chuẩn gồm 10 danh mục: \texttt{00-admin}, \texttt{01-docs}, \texttt{02-repos}, \texttt{03-infra}, \texttt{04-platform}, \texttt{05-ci}, \texttt{06-security}, \texttt{07-observability}, \texttt{08-evidence}, \texttt{09-scripts}, \texttt{10-agent-skill}. Mỗi danh mục có file README mô tả mục đích. \\
\midrule
Clone app-repo &
Đã clone toàn bộ mã nguồn sample app vào \texttt{02-repos/app-repo/}, bao gồm đầy đủ Terraform modules của app (VPC, EKS, tags). \\
\midrule
GitOps repo skeleton &
Đã tạo cấu trúc thư mục cho GitOps Repo theo mô hình Kustomize: \texttt{base/} + \texttt{overlays/dev/}, \texttt{overlays/staging/}, \texttt{overlays/prod/} cùng với thư mục cấu hình Argo CD (\texttt{argocd/applications}, \texttt{argocd/projects}). \\
\midrule
Tài liệu hướng dẫn &
Đã viết \texttt{PROJECT\_ROADMAP.md} (lộ trình 4 phase), \texttt{HUONG\_DAN\_THUC\_HIEN\_CHI\_TIET.md} (hướng dẫn từng bước), và \texttt{START\_HERE\_PHASE1.md} (khởi động nhanh). \\
\midrule
Công cụ local &
Đã cài đặt và xác minh toolchain: Terraform v1.14.8, kubectl v1.34.1, Helm v4.1.3, Docker v29.2.1, AWS CLI v2.31.22. Có script tự động cài đặt (\texttt{setup-tools-windows.ps1}) trong \texttt{09-scripts/}. \\
\bottomrule
\end{longtable}

\subsection{Giai đoạn 2: Infrastructure as Code}

\subsubsection{Terraform Backend — S3 + DynamoDB}

Terraform backend là thành phần nền móng của toàn bộ kiến trúc IaC, đảm bảo:

\begin{itemize}[leftmargin=*,nosep]
  \item \textbf{State lưu tập trung trên S3} — tránh mất mát khi làm việc trên nhiều máy, hỗ trợ collaboration.
  \item \textbf{State locking qua DynamoDB} — ngăn chặn nhiều người cùng apply gây race condition.
  \item \textbf{Versioning + Encryption} — cho phép rollback state và bảo vệ dữ liệu nhạy cảm.
  \item \textbf{Block public access} — bucket không thể truy cập từ internet.
\end{itemize}

\begin{table}[H]
\centering
\begin{tabular}{@{}ll@{}}
\toprule
\textbf{Resource}         & \textbf{Chi tiết} \\
\midrule
S3 Bucket                 & \texttt{uit-devsecops-gitops-dev-0a502713-tfstate} \\
DynamoDB Table            & \texttt{uit-devsecops-gitops-dev-0a502713-tflock} \\
Mã hóa                    & AES256 (server-side encryption) \\
Phiên bản                 & Enabled \\
Public access             & Block tất cả \\
Terraform module          & Tự viết (72 dòng), không dùng module bên ngoài \\
\bottomrule
\end{tabular}
\end{table}

\begin{lstlisting}[caption={Cấu hình backend cho môi trường dev},label={lst:backend}]
# backend.hcl
bucket         = "uit-devsecops-gitops-dev-0a502713-tfstate"
key            = "envs/dev/terraform.tfstate"
region         = "ap-southeast-1"
dynamodb_table = "uit-devsecops-gitops-dev-0a502713-tflock"
encrypt        = true
\end{lstlisting}

\subsubsection{VPC — Mạng nền tảng}

VPC được tạo bằng Terraform module \texttt{terraform-aws-modules/vpc/aws} v5.21.0 với cấu hình tối ưu cho chi phí:

\begin{table}[H]
\centering
\begin{tabular}{@{}ll@{}}
\toprule
\textbf{Thông số}         & \textbf{Giá trị} \\
\midrule
VPC ID                    & \texttt{vpc-0d37b4e9b21e95421} \\
CIDR Block                & \texttt{10.0.0.0/16} \\
Public Subnets            & 3 (mỗi /24, 1 mỗi Availability Zone) \\
Private Subnets           & 3 (mỗi /24, 1 mỗi AZ) \\
NAT Gateway               & 1 (single, thay vì 3 để tiết kiệm) \\
Internet Gateway          & Đã tạo và attach \\
DNS Hostnames             & Enabled \\
\end{tabular}
\end{table}

Các subnet được gán tag tự động cho EKS nhận diện:

\begin{itemize}[leftmargin=*,nosep]
  \item Public subnet: \texttt{kubernetes.io/role/elb = 1} — cho phép tạo Internet-facing Load Balancer.
  \item Private subnet: \texttt{kubernetes.io/role/internal-elb = 1} — cho phép tạo Internal Load Balancer.
  \item Tất cả subnet: \texttt{kubernetes.io/cluster/<name> = shared} — EKS biết subnet này thuộc cluster nào.
\end{itemize}

\subsubsection{EKS Cluster — Kubernetes Control Plane}

EKS cluster được tạo bằng module \texttt{terraform-aws-modules/eks/aws} \textasciitilde\ 19.9 với các cấu hình bảo mật:

\begin{table}[H]
\centering
\begin{tabular}{@{}ll@{}}
\toprule
\textbf{Thông số}              & \textbf{Giá trị} \\
\midrule
Tên cluster                     & \texttt{uit-devsecops-dev} \\
Phiên bản Kubernetes            & \texttt{v1.33.11} \\
Control plane endpoint          & Public + Private \\
IAM Role (cluster)              & Đã tạo với trust policy cho \texttt{eks.amazonaws.com} \\
IAM Policies gán cho role       & AmazonEKSClusterPolicy, AmazonEKSVPCResourceController \\
KMS Encryption                  & Đã bật (mã hóa Kubernetes secrets) \\
OIDC Provider                   & Đã tạo (hỗ trợ IRSA) \\
CloudWatch Logs                 & API, Audit, Authenticator → retention 90 ngày \\
\end{tabular}
\end{table}

\begin{lstlisting}[caption={Khai báo EKS cluster trong Terraform},label={lst:eks}]
module "eks_cluster" {
  source  = "terraform-aws-modules/eks/aws"
  version = "~> 19.9"

  cluster_name                   = var.environment_name
  cluster_version                = "1.33"
  cluster_endpoint_public_access = true

  vpc_id     = var.vpc_id
  subnet_ids = var.subnet_ids

  cluster_addons = {
    vpc-cni = {
      before_compute = true
      most_recent    = true
    }
  }
}
\end{lstlisting}

\subsubsection{Managed Node Groups}

Hai managed node group được triển khai trên hai Availability Zone khác nhau để đảm bảo high availability:

\begin{table}[H]
\centering
\begin{tabular}{@{}lll@{}}
\toprule
\textbf{Node Group}          & \textbf{Instance} & \textbf{Scaling (Min/Desired/Max)} \\
\midrule
\texttt{managed-nodegroup-1} & t3.medium         & 0 / 0 / 1 \\
\texttt{managed-nodegroup-2} & t3.medium         & 0 / 0 / 1 \\
\bottomrule
\end{tabular}
\end{table}

\begin{itemize}[leftmargin=*,nosep]
  \item Mỗi node group nằm trong một private subnet riêng (AZ riêng).
  \item Instance t3.medium: 2 vCPU, 4GB RAM — đủ chạy 5 microservice của sample app.
  \item Hiện tại node đã được scale về 0 để tiết kiệm chi phí. Khi cần làm việc, chỉ cần thay đổi \texttt{node\_min\_size = 1}, \texttt{node\_desired\_size = 1} và chạy \texttt{terraform apply}.
\end{itemize}

\subsubsection{Kubernetes Addons}

\begin{longtable}{p{0.22\textwidth} p{0.28\textwidth} p{0.46\textwidth}}
\toprule
\textbf{Addon} & \textbf{Namespace / Số pods} & \textbf{Mô tả} \\
\midrule
VPC-CNI &
managed, 2 pods &
Quản lý địa chỉ IP cho pods thông qua Elastic Network Interface (ENI). Cho phép mỗi pod có IP riêng trong VPC. \\
\midrule
CoreDNS &
\texttt{kube-system}, 2 pods &
DNS server nội bộ cho cluster, phân giải tên service thành IP. \\
\midrule
kube-proxy &
\texttt{kube-system}, 2 pods &
Network proxy, quản lý Kubernetes Service (iptables/IPVS rules). \\
\midrule
AWS Load Balancer Controller &
\texttt{kube-system}, 2 pods &
Tự động tạo Application Load Balancer (ALB) hoặc Network Load Balancer (NLB) khi tạo Kubernetes Ingress hoặc Service type LoadBalancer. Đây là thành phần quan trọng để expose service ra internet. \\
\midrule
Cert Manager &
\texttt{cert-manager}, 3 pods &
Tự động cấp và gia hạn chứng chỉ TLS từ Let's Encrypt hoặc các CA khác. Gồm controller, cainjector và webhook. \\
\midrule
OpenTelemetry Operator &
\texttt{opentelemetry-operator-system}, 1 pod &
Operator quản lý OpenTelemetry collector. Hiện collector đang ở chế độ disable, sẽ kích hoạt trong Phase 4 (Observability). \\
\bottomrule
\end{longtable}

Tất cả addons được cài đặt thông qua Terraform (Helm provider cho LB Controller và Cert Manager, \texttt{aws\_eks\_addon} cho VPC-CNI và ADOT), đảm bảo tính nhất quán và khả năng tái tạo.

\section{Kết quả kiểm thử}

Sau khi provision hoàn tất, đã thực hiện các bước xác minh hệ thống:

\subsection{Kết nối cluster}

\begin{lstlisting}[caption={Xác minh kết nối đến EKS}]
$ kubectl cluster-info
Kubernetes control plane is running at
  https://F358B0B27A319352351851239024C3F2.gr7.ap-southeast-1.eks.amazonaws.com
CoreDNS is running at
  https://F358B0B27A319352351851239024C3F2.gr7.../api/v1/.../kube-dns:dns/proxy
\end{lstlisting}

\subsection{Trạng thái worker nodes}

\begin{lstlisting}[caption={Danh sách nodes — tất cả Ready}]
$ kubectl get nodes -o wide
NAME                    STATUS   ROLES    AGE    VERSION              INTERNAL-IP
ip-10-0-10-249...       Ready    <none>   2m     v1.33.11-eks-4136f65 10.0.10.249
ip-10-0-11-155...       Ready    <none>   2m     v1.33.11-eks-4136f65 10.0.11.155
\end{lstlisting}

Cả hai node đều ở trạng thái \texttt{Ready}, chạy Amazon Linux 2023, container runtime containerd, địa chỉ IP nằm trong dải private subnet.

\subsection{Trạng thái tất cả pods}

\begin{lstlisting}[caption={Toàn bộ pods đều Running, không restart}]
$ kubectl get pods -A
NAMESPACE     NAME                               READY   STATUS    RESTARTS
cert-manager  cert-manager-6c4645d66c-vn7th      1/1     Running   0
cert-manager  cert-manager-cainjector-...        1/1     Running   0
cert-manager  cert-manager-webhook-...           1/1     Running   0
kube-system   aws-load-balancer-controller-...    1/1     Running   0
kube-system   aws-load-balancer-controller-...    1/1     Running   0
kube-system   aws-node-s9glr                     2/2     Running   0
kube-system   aws-node-whlpn                     2/2     Running   0
kube-system   coredns-6cfd9cdcf-bmsc6            1/1     Running   0
kube-system   coredns-6cfd9cdcf-xbqv2            1/1     Running   0
kube-system   kube-proxy-5vcwh                   1/1     Running   0
kube-system   kube-proxy-pl9vf                   1/1     Running   0
opentelemetry opentelemetry-operator-...         2/2     Running   0
\end{lstlisting}

Tất cả 12 pods đều ở trạng thái \texttt{Running}, không có pod nào bị restart — cluster hoạt động ổn định.

\subsection{Kiểm tra trên AWS Console}

\begin{itemize}[leftmargin=*,nosep]
  \item \textbf{EKS Dashboard:} cluster \texttt{uit-devsecops-dev} hiển thị trạng thái \texttt{Active}, 2 node groups hiển thị \texttt{Active}.
  \item \textbf{EC2 Dashboard:} 2 instances t3.medium đang chạy trong private subnet.
  \item \textbf{VPC Dashboard:} VPC với đầy đủ 6 subnets, Internet Gateway, NAT Gateway, route tables được cấu hình đúng.
  \item \textbf{CloudWatch:} Log group \texttt{/aws/eks/uit-devsecops-dev/cluster} đã được tạo và đang nhận log từ control plane.
\end{itemize}

\section{Quản lý chi phí}

\subsection{Phân tích chi phí dự kiến}

\begin{table}[H]
\centering
\begin{tabular}{@{}lrr@{}}
\toprule
\textbf{Tài nguyên}                    & \textbf{Đơn giá}     & \textbf{Tháng (ước tính)} \\
\midrule
EKS Control Plane                      & \$0.10/giờ           & \~\$72.00 \\
2 $\times$ EC2 t3.medium (khi chạy)    & \$0.0528/giờ/node    & \~\$76.00 \\
NAT Gateway                            & \$0.045/giờ          & \~\$32.00 \\
\midrule
\textbf{Tổng (full-time)}              &                      & \textbf{\~\$180.00} \\
\textbf{Tổng (node off)}               &                      & \textbf{\~\$104.00} \\
\bottomrule
\end{tabular}
\end{table}

\subsection{Biện pháp tiết kiệm đã triển khai}

\begin{enumerate}[leftmargin=*,nosep]
  \item \textbf{Node groups scale về 0} ngoài giờ làm việc — giảm \~\$76/tháng. Cấu hình đơn giản qua \texttt{terraform.tfvars}: \texttt{node\_min\_size = 0}, \texttt{node\_desired\_size = 0}.
  \item \textbf{Single NAT Gateway} thay vì 3 (mỗi AZ một cái) — giảm chi phí NAT xuống còn 1/3.
  \item \textbf{Không dùng Istio} service mesh (\texttt{istio\_enabled = false}) — tránh overhead tài nguyên và chi phí.
  \item \textbf{Không kích hoạt OpenTelemetry collector} — chỉ cài operator, chưa bật data pipeline.
  \item \textbf{Dùng t3.medium} thay vì t3.large hoặc các instance lớn hơn.
  \item \textbf{Terraform destroy} toàn bộ khi nghỉ dài ngày — state lưu trên S3 nên có thể tái tạo lại chính xác trong \textasciitilde 13 phút.
\end{enumerate}

\section{Kế hoạch tiếp theo}

\begin{table}[H]
\centering
\begin{tabular}{@{}p{0.05\textwidth} p{0.55\textwidth} p{0.1\textwidth} p{0.18\textwidth}@{}}
\toprule
\textbf{\#} & \textbf{Công việc} & \textbf{Giai đoạn} & \textbf{Mức ưu tiên} \\
\midrule
1  & Tạo ECR repositories cho 5 microservice                           & Phase 2 & Cao \\
2  & Viết Jenkinsfile: build $\rightarrow$ test $\rightarrow$ scan $\rightarrow$ push ECR & Phase 3 & Cao \\
3  & Tích hợp SonarQube SAST vào CI pipeline                           & Phase 3 & Cao \\
4  & Tích hợp OWASP Dependency Check                                    & Phase 3 & Trung bình \\
5  & Tích hợp Trivy Container Image Scan                                & Phase 3 & Trung bình \\
6  & Viết Kustomize manifests (base + dev overlay)                      & Phase 4 & Cao \\
7  & Cài đặt và cấu hình Argo CD trên EKS                              & Phase 4 & Cao \\
8  & Tự động cập nhật image tag vào GitOps repo từ CI                   & Phase 4 & Trung bình \\
9  & Demo deploy toàn bộ 5 service lên EKS qua Argo CD                  & Phase 4 & Cao \\
10 & Demo rollback bằng \texttt{git revert} trên GitOps repo             & Phase 4 & Cao \\
11 & Thiết lập CloudWatch + Prometheus + Grafana dashboards             & Phase 4 & Thấp \\
12 & Viết runbooks và cấu hình alerts                                   & Phase 4 & Thấp \\
13 & Tổng hợp evidence, hoàn thiện báo cáo cuối kỳ                      & Tổng kết & Cao \\
\bottomrule
\end{tabular}
\end{table}

\section{Rủi ro và biện pháp kiểm soát}

\begin{longtable}{p{0.22\textwidth} p{0.18\textwidth} p{0.56\textwidth}}
\toprule
\textbf{Rủi ro}               & \textbf{Mức độ}   & \textbf{Biện pháp} \\
\midrule
IAM permissions hạn chế &
Đã xử lý &
IAM user \texttt{uit-devsecops-gitops-deployer} đã được gán quyền AdministratorAccess. Với tài khoản thật cần tuân thủ least-privilege — tài liệu tham khảo policy đã được chuẩn bị. \\
\midrule
Chi phí AWS vượt ngân sách &
Đã kiểm soát &
Scale node về 0 ngoài giờ, single NAT Gateway, dùng t3.medium. Có thể \texttt{terraform destroy} toàn bộ khi nghỉ dài ngày. \\
\midrule
Thời gian provision EKS lâu (\textasciitilde 15 phút) &
Thấp &
Terraform tự động hóa hoàn toàn, không cần can thiệp thủ công. Có thể chạy nền trong lúc làm việc khác. \\
\midrule
Phiên bản module không tương thích &
Thấp &
Tất cả module được pin version cụ thể (\texttt{\textasciitilde\ 19.9}, \texttt{5.21.0}), không dùng \texttt{latest}. \\
\midrule
Quên tắt tài nguyên gây lãng phí &
Đã kiểm soát &
File \texttt{terraform.tfvars} đã được cấu hình sẵn cho chế độ tiết kiệm (node = 0). Khi cần làm việc tiếp chỉ cần sửa lại \texttt{node\_desired\_size = 1} và \texttt{terraform apply}. \\
\bottomrule
\end{longtable}

\section{Cấu trúc mã nguồn}

\begin{verbatim}
do-an-devsecops-gitops/
  00-admin/              Quan ly du an (meeting notes, decision log)
  01-docs/               Tai lieu (report, slides, thesis plan)
    report/              Bao cao tien do
  02-repos/              Ma nguon ung dung & GitOps manifests
    app-repo/            Sample app (aws-containers/retail-store-sample-app)
      src/               5 microservices (ui, cart, orders, catalog, checkout)
      terraform/         Terraform modules cua app (VPC, EKS, tags)
    gitops-repo/         GitOps deployment manifests (Kustomize)
      apps/retail-store/
        base/
        overlays/{dev,staging,prod}/
  03-infra/terraform/    Ha tang cua do an (Terraform)
    backend/             S3 + DynamoDB state backend
    envs/dev/            Moi truong dev (VPC + EKS)
  04-platform/           Kubernetes platform add-ons (Argo CD)
  05-ci/                 Jenkins CI pipeline (Jenkinsfile)
  06-security/           Security artifacts (policies, SBOM, scan reports)
  07-observability/      Monitoring (alerts, dashboards, runbooks)
  08-evidence/           Bang chung cho tung phase
  09-scripts/            Script ho tro (setup toolchain, switch profile)
  10-agent-skill/        AI assistant skill definition
\end{verbatim}

\section{Kết luận}

Đồ án đã hoàn thành giai đoạn lập kế hoạch (Phase 1) và đang tiến hành giai đoạn hạ tầng (Phase 2). Các kết quả chính đạt được đến thời điểm hiện tại:

\begin{enumerate}[leftmargin=*,nosep]
  \item \textbf{Kiến trúc rõ ràng:} Mô hình DevSecOps + GitOps đã được thiết kế chi tiết, bao gồm sơ đồ luồng dữ liệu, danh sách AWS services, và kế hoạch triển khai theo 4 phase.
  \item \textbf{Hạ tầng sẵn sàng:} VPC, EKS cluster v1.33, và toàn bộ addons (Load Balancer Controller, Cert Manager, OTEL Operator) đã được provision thành công bằng Terraform, kiểm thử và xác minh hoạt động ổn định.
  \item \textbf{Chi phí được kiểm soát:} Đã triển khai các biện pháp tiết kiệm (scale node về 0, single NAT, instance t3.medium), tổng chi phí tối đa \~\$104/tháng, có thể giảm về \$0 khi \texttt{terraform destroy}.
  \item \textbf{Tài liệu đầy đủ:} Workspace được tổ chức khoa học, có sẵn tài liệu roadmap, hướng dẫn thực hiện, và script tự động hóa.
\end{enumerate}

Các bước tiếp theo sẽ tập trung vào việc xây dựng CI pipeline với security gates (Jenkins + SonarQube + OWASP + Trivy), tạo Kustomize manifests, và triển khai CD qua Argo CD.

\end{document}
